% IEEE Control Systems Letters template
\documentclass[journal]{IEEEtran}

%--------------- Packages ---------------%
\usepackage[utf8]{inputenc}
\usepackage{amsmath,amsfonts,amssymb,amsthm}
\usepackage{graphicx}
\usepackage{cite}
\usepackage{bm}
\usepackage{siunitx}
\usepackage{url}
\usepackage{color}
\usepackage{mathtools}
\usepackage{mathrsfs}

\newtheorem*{problem}{Problem Statement}
\newtheorem{definition}{Definition}
\newtheorem{proposition}{Proposition}

%--------------- Macros ---------------%
\newcommand{\norm}[1]{\left| #1 \right|}
\newcommand{\tp}{^\top}
\newcommand{\ball}{\mathbb{B}}
\newcommand{\eye}[1]{I_{#1}}
\DeclareMathOperator{\trace}{trace}
\DeclareMathOperator{\crit}{crit}

%--------------- Title & Authors ---------------%
\title{CubeSat Drone: The Development of a CubeSat Experimental Testbed}
% Trajectory Tracking and Control Allocation for an Overactuated 3U CubeSat with Ten Propellers

\author{Miguel~A.~Nunes, Pedro~Casau,~\IEEEmembership{Member,~IEEE}% <-this % stops a space
\thanks{Miguel A. Nunes is with the Hawai‘i Space Flight Laboratory at University of Hawai‘i, Honolulu, USA. P. Casau is with the Department of Electronics, Telecommunications and Informatics at University of Aveiro, and with Instituto de Telecomunicações at Aveiro, Portugal. Corresponding author: Miguel A. Nunes (e-mail: \texttt{miguel.nunes@hsfl.hawaii.edu}).}% <-this % stops a space
}

%--------------- Begin Document ---------------%
\begin{document}
\maketitle

%--------------- Abstract ---------------%
\begin{abstract}
This paper presents a hybrid control architecture for global trajectory tracking of an aerial vehicle with the form factor of a \qty{3}{U} CubeSat. 
The vehicle is equipped with ten bidirectional propellers mounted on facets of its frame. 
We derive a dynamical model that accounts for the propeller layout and design a proportional-derivative (PD) feedback law combined with a hybrid supervisor to globally asymptotically stabilize the vehicle's pose along feasible reference trajectories. 
To solve the control allocation problem, we implement a minimum control effort algorithm.
Simulation results demonstrate convergence of position and orientation errors and produce physically plausible rotor speed commands.
\end{abstract}

\begin{IEEEkeywords}
Overactuated systems, control allocation, CubeSat, multirotor, hybrid systems simulation
\end{IEEEkeywords}

%--------------- Introduction ---------------%
\section{Introduction}
\IEEEPARstart{M}{ultirotor} spacecraft employing electric propulsion have recently been proposed to expand the agility of small satellites during proximity operations and planetary exploration. Compared to conventional thruster arrangements, arrays of miniature propellers (ducted fans) can provide continuous, fully bidirectional control forces while fitting within the \qtyproduct{10 x 20 x 10}{cm} envelope of a \qty{2}{U} CubeSat. However, when more actuators than controlled degrees of freedom are available, the designer must address \\emph{control allocation}---the mapping of desired net forces and moments to individual actuator commands.

This letter considers a CubeSat carrying ten fixed--pitch propellers, one centered on each outward--facing facet (Fig.~\ref{fig:layout}). The arrangement grants a rich actuation set but introduces redundancy. Building on classical rigid--body dynamics and linear algebraic allocation, we propose a PD controller that demands only onboard measurements of position, velocity, attitude, and angular rate. The method inverts, at each control step, an analytically derived allocation matrix and thus avoids iterative optimization.

Our contribution is twofold: (i)~we formulate the full six--degree--of--freedom dynamics of the overactuated CubeSat in compact matrix form, and (ii)~we validate a lightweight control-allocation scheme in numerical simulation, highlighting its ability to respect actuator limits while achieving tight trajectory tracking.

%--------------- Problem Statement ---------------%
\section{Problem Statement}
\subsection{Notation}
The 3-dimensional ball of radius $L\geq 0$ is represented by $L\ball = \{\bm{x}\in\mathbb{R}^3:\norm{\bm{x}}\leq L\}$. The special orthogonal group is $\mathrm{SO}(3) = \{R\in\mathbb{R}^{3\times3}:RR\tp=\eye{3}, \det(R)=1\}$. The identity matrix is $\eye{n}\in\mathbb{R}^{n\times n}$. 
The function 
\begin{equation}\label{eq:S}
S(\omega)\coloneqq\begin{bmatrix}
0 & -\omega_3 & \omega_2\\
\omega_3 & 0 & -\omega_1\\
-\omega_2 & \omega_1 & 0
\end{bmatrix}
\end{equation}
maps $\bm{\omega}\coloneqq(\omega_1,\omega_2,\omega_3)\in\mathbb{R}^3$ to a skew-symmetric matrix such that $S(\bm{\omega})\bm{x} = \bm{\omega}\times\bm{x}$ for any $\bm{x}\in\mathbb{R}^3$.

\begin{definition}[Potential Function on $SO(3)$]
A function $U:SO(3)\to\mathbb{R}$ is said to be a potential function on $SO(3)$ if it is continuously differentiable and positive definite relative to $\eye{3}$.
\end{definition}

\begin{definition}
The set of critical points of a potential function $U:SO(3)\to\mathbb{R}$ is defined as
$$\crit U = \{R\in SO(3): \nabla U(R) = 0\}.$$
\end{definition}

\begin{definition}
Given a discrete set $\mathcal{Q}$, a collection of potential functions $\mathscr{U_q}\coloneqq\{U_q\}_{q\in\mathcal{Q}}$ on $SO(3)$ is said to be a synergistic if, for each $U_q\in\mathscr{U}$ and each $R\in\crit U_q\setminus\{\eye{3}\}$, there exists another potential function $U_h$ in the collection such that $U_q(R) - U_h(R)>0$.
\end{definition}

\subsection{Rigid-Body Dynamics}
Let $\bm{p}\in\mathbb{R}^3$ denote the position of the vehicle's center of mass in an inertial frame $\{\mathcal{I}\}$, $\bm{v}=\dot{\bm{p}}$ its linear velocity, $R\in\mathrm{SO}(3)$ the body-to-inertial rotation matrix, and $\bm{\omega}\in\mathbb{R}^3$ the angular velocity expressed in the body frame.
The translational and rotational equations of motion are
\begin{equation}\label{eq:dynamics}
\begin{aligned}
\dot{\bm{p}} &= \bm{v},\\
 m\dot{\bm{v}} &= m\bm{g} + R\bm{F},\\
 \dot{R} &= RS(\bm{\omega}),\\
 J\dot{\bm{\omega}} &= \bm{M} - \bm{\omega}\times J\bm{\omega},
\end{aligned}
\end{equation}
where $m$ is the mass, $\bm{g}$ is the gravity vector, $J\in\mathbb{R}^{3\times3}$ is the tensor of inertia, $\bm{F}\in\mathbb{R}^3$ is the net force expressed in the body frame, and $\bm{M}\in\mathbb{R}^3$ is the net moment (see e.g.~\cite{bullo_geometric_2005}).

\subsection{Actuation Model}
Ten identical propellers are located at positions $\bm{r}_i\in\mathbb{R}^3$ and produce thrust along unit vectors $\bm{u}_i$. The $i$-th propeller's rotation speed $\Omega_i\in\mathbb{R}$ generates a force
$$f_i = \rho\, d^4\, C_T\,\norm{\Omega_i}\Omega_i$$ 
where $\rho$ is atmospheric density, $d$ is the propeller's diameter, and $C_T$ is the coefficient of thrust. 
Each propeller also induces a moment about its axis of rotation
$$m_i = \rho d^5 C_Q\, d_i\,\norm{\Omega_i}\Omega_i,$$
where $C_Q$ is the coefficient of torque, and $d_i\in\{+1,-1\}$ captures the direction of rotation (cf.~\cite{VonMises1959}).

The net force and moment acting on the vehicle are
\begin{equation}\label{eq:FM}
\begin{aligned}
\bm{F} &= \sum_{i=1}^{10} f_i\bm{u}_i = k_f\sum_{i=1}^{10}\norm{\Omega_i}\Omega_i\bm{u}_i,\\
\bm{M} &= \sum_{i=1}^{10}\bigl(\bm{r}_i\times f_i\bm{u}_i + m_i\bm{u}_i\bigr) \\
&= \sum_{i=1}^{10}\bigl(k_fS(\bm{r}_i) - k_m d_i\bigr)\bm{u}_i\norm{\Omega_i}\Omega_i,
\end{aligned}
\end{equation}
where $k_f = \rho d^4 C_T$ and $k_m = \rho d^5 C_Q$. Defining $\bm{\Omega}\coloneqq(\norm{\Omega_1}\Omega_1,\ldots,\norm{\Omega_{10}}\Omega_{10})\in\mathbb{R}^{10}$ as the vector of squared propeller speeds, we can rewrite~\eqref{eq:FM} as $\begin{bmatrix} \bm{F}\tp & \bm{M}\tp \end{bmatrix}\tp = T\bm{\Omega}$ with 
\begin{equation*}
T = \begin{bmatrix}
 k_f\bm{u}_1 &\ldots& k_f\bm{u}_{10}\\
 \left(k_fS(\bm{r}_1) - k_m d_1\right)\bm{u}_1 &\ldots& \left(k_fS(\bm{r}_{10}) - k_m d_{10}\right)\bm{u}_{10}
 \end{bmatrix}.
\end{equation*}

\subsection{Control Objective}
Given $L_v,L_\omega>0$ and a twice-differentiable reference trajectory $t\mapsto(\bm{p}_d(t),R_d(t))\in\mathbb{R}^3\times SO(3)$ with bounded derivatives satisfying
\begin{equation}\label{eq:ref}
\begin{aligned}
  \dot{\bm{p}}_d(t) &= \bm{v}_d(t), &
  \dot{\bm{v}}_d(t) &\in L_v\ball\\
  \dot{R}_d(t) &= R_d(t)S(\bm{\omega}_d(t)),& 
  \dot{\bm{\omega}}_d(t) &\in L_\omega\ball
\end{aligned}
\end{equation}
for each $t\geq 0$, we define the tracking errors as follows:
\begin{equation}\label{eq:errors}
\begin{aligned}
\bm{p}_e(t)&\coloneqq\bm{p}(t)-\bm{p}_d(t),&
\bm{v}_e(t)&\coloneqq\bm{v}(t)-\dot{\bm{p}}_d(t),\\
R_e(t)&\coloneqq R_d(t)^\top R(t)&
\bm{\omega}_e(t) &\coloneqq \bm{\omega}(t)-R_e(t)\tp\bm{\omega}_d(t)
\end{aligned}
\end{equation}
for each $t\geq 0$. 
Differentiating~\eqref{eq:errors} with respect to time and using~\eqref{eq:dynamics} and~\eqref{eq:ref}, we obtain the error dynamics
\begin{equation}\label{eq:errordyn}
\begin{aligned}
\dot{\bm{p}}_e &= \bm{v}_e\\
m\dot{\bm{v}}_e &= m\bm{g} + R_d R_e \bm{F} - m\ddot{\bm{p}}_d,\\
\dot{R}_e &= R_e S(\bm{\omega}_e),\\
J\dot{\bm{\omega}_e} &= \Sigma(R_e,\bm{\omega}_e,\bm{\omega}_d)\bm{\omega}_e+\bm{M} + S(JR_e\tp\bm{\omega}_d)R_e\tp\bm{\omega}_d - R_e\tp\dot{\bm{\omega}}_d
\end{aligned}
\end{equation}
with
\begin{multline*}
\Sigma(R_e,\bm{\omega}_e,\bm{\omega}_d) = S(J\bm{\omega}_e)+S(JR_e\tp\bm{\omega}_d)\\
-S(R_e\tp\bm{\omega}_d)J-J S(R_e\tp\bm{\omega}_d),
\end{multline*}
which is a skew-symmetric matrix. 

\begin{problem}
  Design a controller that globally asymptotically stabilizes the identity for~\eqref{eq:errordyn}.
\end{problem}

%--------------- Controller Design ---------------%
\section{Controller Design}

The controller design followed in this paper extends the control strategy proposed in~\cite{mayhew_synergistic_2013} from attitude tracking to full pose tracking. The proposed controller is based on rotation matrix feedback combined with a minimum effort control allocation strategy. Rotation matrix feedback is preferred over other attitude representations, such as Euler angles or unit-quaternions, because it is free of singularities or ambiguities. In addition,~\cite{mayhew_synergistic_2013} provides us with a simple way of achieving global asymptotic stability through hybrid control, which, for the sake of completeness, is explained in the following sections.

\subsection{Nominal Feedback Law Design}
Almost global attitude tracking is achieved through the appropriate design of a Lyapunov function and feedback pair. In this direction, we consider the following candidate Lyapunov function:
\begin{equation*}
  V(R_e,\bm{\omega}_e) = U_0(R_e)+\frac{1}{2}\bm{\omega}_e\tp J\bm{\omega}_e,
\end{equation*}
where 
\begin{equation}\label{eq:U0}
  U_0(R_e) = \trace(A(\eye{3}-R_e))
\end{equation}
is a modified trace function with $A\in\mathbb{R}^{3\times3}$ positive definite. The time-derivative of $V$ along the trajectories of~\eqref{eq:errordyn} is
\begin{multline*}
  \dot{V}(R_e,\bm{\omega}_e) = -\bm{\omega_e}\tp S^{-1}(R_e\tp A - A R_e) \\
  +\bm{\omega}_e\tp(\Sigma(R_e,\bm{\omega}_e,\bm{\omega}_d)\bm{\omega}_e+\bm{M} + S(JR_e\tp\bm{\omega}_d)R_e\tp\bm{\omega}_d - R_e\tp\dot{\bm{\omega}}_d)
\end{multline*}
where $S^{-1}$ is the inverse of the skew-symmetric map~\eqref{eq:S}. Defining the feedback law as 
\begin{equation}\label{eq:attctrl}
  \bm{M} = S^{-1}(R_e\tp A - A R_e) - K_\omega \bm{\omega}_e - S(JR_e\tp\bm{\omega}_d)R_e\tp\bm{\omega}_d + R_e\tp\dot{\bm{\omega}}_d,
\end{equation}
with $K_\omega\in\mathbb{R}^{3\times3}$ positive definite, we obtain
\begin{equation*}
  \dot{V}(R_e,\bm{\omega}_e) = -\bm{\omega}_e\tp K_\omega \bm{\omega}_e \leq 0.
\end{equation*}

Stabilization for the tracking error dynamics of the position subsystem is achieved through a standard PD controller:
\begin{equation}\label{eq:Fctrl}
\bm{F} = -R_e\tp R_d\tp(K_p\bm{p}_e - K_v\bm{v}_e - m\bm{g})+m\ddot{\bm{p}}_d,
\end{equation}
with positive definite gains $K_p,K_v\in\mathbb{R}^{3\times 3}$.

\subsection{Complementary Feedback Law Design}
The feedback laws~\eqref{eq:attctrl} and~\eqref{eq:Fctrl} do not guarantee global asymptotic stability due to the topological obstruction of $\mathrm{SO}(3)$ (cf.~\cite{bhat_topological_2000}). In particular, $U_0$ has four critical points, one corresponding to the desired equilibrium $R_e=\eye{3}$ and three undesired equilibria corresponding to rotations of $180^{\circ}$ about the eigenvectors of $A$ (cf.~\cite{mayhew_synergistic_2013}). To overcome this limitation, a complementary feedback law is designed by selecting a different modified trace function $U_1$ with the same structure as $U_0$ but with different critical points, namely:
\begin{equation}\label{eq:U1}
  U_1(R_e) = \alpha+\epsilon U_0(P\tp R_e),
\end{equation} 
where $\alpha,\epsilon>0$ and $P\in SO(3)$ are controller parameters. The feedback law associated with $U_1$ is
\begin{multline}\label{eq:attctrl2}
  \bm{M} = \epsilon S^{-1}(R_e\tp P A - A P\tp R_e) - K_\omega \bm{\omega}_e \\
  - S(JR_e\tp\bm{\omega}_d)R_e\tp\bm{\omega}_d + R_e\tp\dot{\bm{\omega}}_d.
\end{multline}
The values of $\alpha,\epsilon$ and $P$ are chosen in combination with the Hybrid Supervisor described in the next section so as to avoid undesired equilibria. The force input remains unchanged relative to~\eqref{eq:Fctrl}.

\subsection{Hybrid Supervisor}
The Hybrid Supervisor proposed in~\cite[Example~9]{mayhew_synergistic_2013} continuously evaluates the difference between the current value of $U_q(R_e)$ and $U_{1-q}(R_e)$ and switches from $q$ to $1-q$ whenever the difference exceeds a threshold $\delta>0$. The switching logic is defined by the hybrid system $\mathcal{H}\coloneqq (C,F,D,G)$ with state $(\xi_d,\xi_e,q)\in\mathcal{S}_d\times\mathcal{S}\times \mathcal{Q}$ with $\xi_{d,e}\coloneqq (\bm{p}_{d,e},\bm{v}_{d,e},R_{d,e},\bm{\omega}_{d,e})$, $\mathcal{S}_{d}=\mathcal{S}_{e}\coloneqq \mathbb{R}^3\times\mathbb{R}^3\times SO(3)\times\mathbb{R}^3$ and $\mathcal{Q}\coloneqq\{0,1\}$, as follows:
\begin{align*}
  \dot{q}&=0 & (\xi_d,\xi_e,q)&\in C\\
  q^+&=1-q & (\xi_d,\xi_e,q)&\in D\\
  & & (F,M)&=(\kappa_F(\xi_d,\xi_e),\kappa_M(\xi_d,\xi_e,q))
\end{align*}
where $\kappa_F$ is defined by~\eqref{eq:Fctrl} and $\kappa_M$ is defined by~\eqref{eq:attctrl} if $q=0$ and by~\eqref{eq:attctrl2} if $q=1$. The flow and jump sets are
\begin{align*}
  C &= \{\xi\in\mathcal{S}: U_q(R_e)-U_{1-q}(R_e)\leq\delta\},\\
  D &= \{\xi\in\mathcal{S}: U_q(R_e)-U_{1-q}(R_e)\geq\delta\}.
\end{align*}

\subsection{Closed-Loop Hybrid System}
The closed-loop hybrid system is given by
\begin{equation}\label{eq:H}
\begin{aligned}
\underbrace{%
\begin{aligned}
&\dot{\bm{p}}_d = \bm{v}_d\\
&\dot{\bm{v}}_d \in L_v\ball\\
&\dot{R}_d = R_d \operatorname{S}(\bm{\omega}_d)\\
&\dot{\bm{\omega}}_d \in L_\omega\ball\\
&\dot{\bm{p}}_e = \bm{v}_e\\
&m\dot{\bm{v}}_e = m\bm{g} + R_d R_e \kappa_F(\xi_d,\xi_e) - m\ddot{\bm{p}}_d\\
&\dot{R}_e = R_e \operatorname{S}(\bm{\omega}_e)\\
&\begin{multlined}
J\dot{\bm{\omega}}_e
= \Sigma(R_e,\bm{\omega}_e,\bm{\omega}_d)\,\bm{\omega}_e + \kappa_M(\xi_d,\xi_e,q) \\
+ \operatorname{S}\!\big(J R_e^\top \bm{\omega}_d\big)\,R_e^\top \bm{\omega}_d
- R_e^\top \dot{\bm{\omega}}_d
\end{multlined}\\
&\dot{q} = 0
\end{aligned}}_{\displaystyle(\xi_d,\xi_e,q)\in C}
\quad
\underbrace{%
\begin{aligned}
\bm{p}_d^{+} &= \bm{p}_d\\
\bm{v}_d^{+} &= \bm{v}_d\\
R_d^{+} &= R_d\\
\bm{\omega}_d^{+} &= \bm{\omega}_d\\
\bm{p}_e^{+} &= \bm{p}_e\\
\bm{v}_e^{+} &= \bm{v}_e\\
R_e^{+} &= R_e\\
\bm{\omega}_e^{+} &= \bm{\omega}_e\\
\\
q^{+} &= 1 - q
\end{aligned}}_{\displaystyle(\xi_d,\xi_e,q)\in D}
\end{aligned}
\end{equation}
to which the following result applies.

\begin{proposition}
  Given $\mathcal{Q}\coloneqq\{0,1\}$, $U_0$ in~\eqref{eq:U0} and $U_1$ in~\eqref{eq:U1}, if $\mathscr{U}\coloneqq\{U_q\}_{q\in\mathcal{Q}}$ is synergistic with gap exceeding $\delta>0$, then the compact set 
  $$\mathcal{A}\coloneqq\{(\xi_d,\xi_e,q)\in\mathcal{S}_d\times\mathcal{S}_e\times\mathcal{Q}:\xi_e=(0,0,\eye{3},0),q=0\}$$
  is globally asymptotically stable for the hybrid system~\eqref{eq:H}.
\end{proposition}

\subsection{Pseudoinverse Allocation}
The control allocation problem consists of finding a vector of propeller speeds 
\begin{equation*}
\bm{\Omega}\coloneqq \begin{bmatrix}
\norm{\Omega_1}\Omega_1 & \ldots & \norm{\Omega_{10}}\Omega_{10}\end{bmatrix}\tp
\end{equation*}
that generates the desired force and moment $(\bm{F},\bm{M})$ given by the feedback laws $\kappa_F$ and $\kappa_M$. We say that the control is minimum effort because it solves the following optimization problem:
\begin{equation}\label{eq:allocation}
\bm{\Omega}=\arg\min\left\{\norm{\bm{\Omega}}^2: T\bm{\Omega} = \begin{bmatrix} \kappa_F(\xi_d,\xi_e)\\ \kappa_M(\xi_d,\xi_e,q) \end{bmatrix}\right\}.
\end{equation}
If $T$ is full rank, then there is a unique solution to~\eqref{eq:allocation} given by
\begin{equation}
\bm{\Omega} = T\tp\left(TT\tp\right)^{-1}\begin{bmatrix} \kappa_F(\xi_d,\xi_e)\\ \kappa_M(\xi_d,\xi_e,q) \end{bmatrix}.
\end{equation}

%--------------- Simulation Results ---------------%
\section{Simulation Results}
Numerical simulations were conducted in MATLAB/Simulink using the \texttt{HyEQsolver} toolbox~\cite{HyEQ2012}. Table~\ref{tab:params} lists the physical parameters adopted from commercial off--the--shelf components.

\begin{table}[!t]
  \caption{Simulation Parameters}
  \label{tab:params}
  \centering
  \begin{tabular}{lc}
    \hline
    Parameter & Value \\
    \hline
    Mass $m$ & \qty{4}{kg} \\
    Inertia $J$ & $\operatorname{diag}(6.7,33.3,33.3)\times10^{-3}\,\si{kg\,m^{2}}$ \\
    Thrust coeff. $k_f$ & \num{1}\,-- (normalized) \\
    Torque coeff. $k_m$ & \num{1}\,-- (normalized) \\
    Gravity $g$ & \qty{9.81}{m\per s^{2}} \\
    Gains $K_p,K_v$ & $\operatorname{diag}(1,1,1)$ \\
    Gains $K_R,K_\omega$ & $\operatorname{diag}(1,1,1)$ \\
    \hline
  \end{tabular}
\end{table}

The CubeSat was required to stabilize from an initial offset of \qty{0.5}{m} along $+x$ and an attitude error of $45^{\circ}$ about the $x$--axis. Figures~\ref{fig:quat} and~\ref{fig:rpm} depict the time histories of the unit quaternion and rotor speeds, respectively. All states converge within \qty{10}{s}, and rotor commands remain within feasible limits. A frame--by--frame animation (not shown) confirms smooth, collision--free motion.

%\begin{figure}[!t]
%  \centering
%  \includegraphics[width=0.95\columnwidth]{quat.pdf}
%  \caption{Quaternion components $q_0,q_1,q_2,q_3$ during the stabilization maneuver.}
%  \label{fig:quat}
%\end{figure}

%\begin{figure}[!t]
%  \centering
%  \includegraphics[width=0.95\columnwidth]{rpm.pdf}
%  \caption{Propeller speeds obtained via the allocation law \eqref{eq:allocation}.}
%  \label{fig:rpm}
%\end{figure}

%--------------- Conclusions ---------------%
\section{Conclusions}
This letter demonstrated that a conventional PD outer loop combined with a pseudoinverse control allocator suffices to govern an overactuated \qty{2}{U} CubeSat outfitted with ten propellers. While simple, the design reliably tracks pose commands without requiring iterative optimization. Future work will address actuator saturation, aerodynamic coupling in rarefied atmospheres, and experimental validation aboard a hardware demonstrator.

%--------------- References ---------------%
\bibliographystyle{ieeetr}
\bibliography{biblio}

\end{document}

